\documentclass[10pt,a4paper]{amsart}
\usepackage[margin=19mm]{geometry}
\usepackage{amsmath,amssymb,mathtools}
\usepackage[scr=rsfs,cal=euler]{mathalfa}
\usepackage{xfrac}
\usepackage[unicode,hidelinks,bookmarks=false]{hyperref}
\newcommand{\ie}{i.e.\ }
\newcommand{\cf}{cf.\ }
\newcommand{\resp}{respectively\ }
\newcommand{\Subsection}{\subsection}
\providecommand{\nobreakdash}{\nobreak\hbox{-}}
\setlength{\emergencystretch}{2em}
\multlinegap0pt
\usepackage{amssymb}
\usepackage{float}
\usepackage[normalem]{ulem}
\usepackage{xcolor}
\newtheorem{lemma}{Lemma}
\newcommand{\smat}[4]{\left(\begin{smallmatrix}#1&#2\\#3&#4\end{smallmatrix}\right)}
\title{Eta-Quotient Representations for a Three Parameter Family of Modular Functions Associated with the Rogers-Ramanujan Continued Fraction}
\author{Bishnu Paudel, James A. Sellers, Haiyang Wang}
\date{Source: arXiv:2609.13131v1 (CC BY 4.0)}
\begin{document}
\maketitle

\noindent\emph{Source and licence.} Eta-Quotient Representations for a Three Parameter Family of Modular Functions Associated with the Rogers-Ramanujan Continued Fraction. Version arXiv:2609.13131v1 (CC BY 4.0), distributed by arXiv under the Creative Commons Attribution 4.0 International licence, http://creativecommons.org/licenses/by/4.0/, as recorded in the arXiv licence metadata for this exact version and shown on the abstract page https://arxiv.org/abs/2609.13131v1. Full licence notice: \texttt{licenses/arxiv-2609.13131-CC-BY-4.0.txt}.\par\smallskip

\noindent\textit{Editorial scope.} Counted unit: the article's lemma lem.transform (source0.tex lines 600--609). The article's complete original proof of it is delivered with it (source0.tex lines 611--714, one \texttt{proof} environment of 104 lines). No mathematics is written, repaired or re-derived: the statement and the proof are the article's own lines, in the article's own notation, and no retained line is substituted or re-flowed.  Retained uncounted as the source-local context the proof needs: lines 516--518; lines 520--522; lines 524--529; lines 542--545; lines 548--551; lines 591--594.  Nothing was omitted from any retained span, so the package carries no omission record.  No retained line contains a citation, so the wrapper carries no bibliography and no bibliography entry.  No retained line contains a figure environment, an image-inclusion command or drawing code, so no image is delivered and the package declares no asset dependency.  Editorial formatting only, in addition: an AMS \texttt{amsart} preamble that restates the article's own declarations and macros used by the retained lines, namely the environments \texttt{lemma} on separate counters, together with the standard packages those lines need.  The retained environments print in this document as Lemma 1 in order of appearance, whereas the article prints the same items with the numbering of its own full document.  The complete native source of the article as distributed in the arXiv e-print for this exact version is bundled unchanged as \texttt{sources/210185/source0.tex}.
\begin{equation}\label{eq.Nshift}
\eta_{N,g+N}(\tau)=\eta_{N,-g}(\tau)=-\eta_{N,g}(\tau).
\end{equation}

\begin{equation}\label{eq.translation}
\eta_{N,g}(\tau+b)=e^{\pi ibNB_{2}(g/N)}\,\eta_{N,g}(\tau).
\end{equation}

\begin{equation}\label{eq.cne0}
\eta_{N,g}(\gamma\tau)
=\varepsilon(a,bN,c/N,d)\,
e^{\pi i\left(\frac{g^{2}ab}{N}-gb\right)}\,
\eta_{N,ag}(\tau),
\end{equation}

\begin{equation}
\eta_{N,g+kN}(\tau)=(-1)^k\eta_{N,g}(\tau),
\label{eq.indexshift}
\end{equation}

\begin{equation}
\eta_{N,N-g}(\tau)=\eta_{N,g}(\tau).
\label{eq.reflect}
\end{equation}

\begin{equation}\label{eq.rgeneta}
r_{j}(\tau)=\frac{\eta_{5j,j}(\tau)}{\eta_{5j,2j}(\tau)}
\qquad(j=1,2,3).
\end{equation}

\begin{lemma}\label{lem.transform}
Let $\gamma=\smat abcd\in\Gamma_{0}(30)$ and $j\in\{1,2,3\}$.  Then
\begin{equation}\label{eq.gamact}
r_{j}(\gamma\tau)=
\begin{cases}
\zeta^{jab}\,r_{j}(\tau), & a\equiv\pm1\pmod5,\\[3pt]
-\,\zeta^{jab}\big/r_{j}(\tau), & a\equiv\pm2\pmod5.
\end{cases}
\end{equation}
\end{lemma}

\begin{proof}
If $c=0$, then $ad=1$, so $a=d=\pm1$.  Hence
\begin{equation*}
\gamma\tau=\frac{a\tau+b}{a}=\tau+ab.
\end{equation*}
Applying \eqref{eq.translation} to the
numerator and denominator in \eqref{eq.rgeneta} gives
\begin{equation*}
\begin{aligned}
r_{j}(\gamma\tau)
&=r_{j}(\tau+ab)\\
&=e^{\pi i\cdot5jab
\left(B_{2}(\frac15)-B_{2}(\frac25)\right)}r_{j}(\tau)\\
&=e^{2\pi ijab/5}r_{j}(\tau)
=\zeta^{jab}r_{j}(\tau),
\end{aligned}
\end{equation*}
which proves \eqref{eq.gamact} when $c=0$.

Now suppose that $c\ne0$. 
Since $5j\mid30\mid c$,  we may apply \eqref{eq.cne0}, with  $N=5j$, to the
numerator and denominator in \eqref{eq.rgeneta}.  The
$\varepsilon$-factors are the same and cancel.  Therefore

\begin{equation}\label{eq.transformraw}
\begin{aligned}
r_{j}(\gamma\tau)
&=
e^{\pi i\left(
\frac{(j^{2}-4j^{2})ab}{5j}-(j-2j)b
\right)}
\frac{\eta_{5j,aj}(\tau)}
{\eta_{5j,2aj}(\tau)}\\
&=
e^{\pi i\,jb(5-3a)/5}
\frac{\eta_{5j,aj}(\tau)}
{\eta_{5j,2aj}(\tau)}.
\end{aligned}
\end{equation}

Because $a$ is odd, $u:=\tfrac12(5-3a)\in\mathbb{Z}$.  As
$2u\equiv2a\pmod5$, we have $u\equiv a\pmod5$, and thus
\begin{equation}
e^{\pi i\,\frac{jb(5-3a)}{5}}=e^{2\pi i\,jbu/5}=\zeta^{jbu}=\zeta^{jab}.
\label{eq.expuniform}
\end{equation}

Since $a$ is odd and coprime to $5$, we may write
\begin{equation*}
a=10t+a_0,
\qquad
a_0\in\{1,9,7,3\},
\end{equation*}
where these four choices correspond, respectively, to
$a\equiv1,-1,2,-2\pmod5$. We then have
\begin{equation*}
aj=a_0j+2t(5j),
\qquad
2aj=2a_0j+4t(5j).
\end{equation*}
Thus, by \eqref{eq.indexshift},
\begin{equation*}
\frac{\eta_{5j,aj}(\tau)}
{\eta_{5j,2aj}(\tau)}
=
\frac{\eta_{5j,a_0j}(\tau)}
{\eta_{5j,2a_0j}(\tau)}.
\end{equation*}	
Using \eqref{eq.Nshift}, \eqref{eq.indexshift} and \eqref{eq.reflect}, we obtain
\begin{equation*}
\begin{array}{c|c|c|c}
a_{0}
& a\pmod5
& \eta_{5j,a_{0}j}(\tau)
& \eta_{5j,2a_{0}j}(\tau)\\ \hline
1 & 1
& \eta_{5j,j}(\tau)
& \eta_{5j,2j}(\tau)\\
9 & -1
& -\eta_{5j,j}(\tau)
& -\eta_{5j,2j}(\tau)\\
7 & 2
& -\eta_{5j,2j}(\tau)
& \eta_{5j,j}(\tau)\\
3 & -2
& \eta_{5j,2j}(\tau)
& -\eta_{5j,j}(\tau).
\end{array}
\end{equation*}
Therefore,
\begin{equation}\label{eq.quot}
\frac{\eta_{5j,aj}(\tau)}
{\eta_{5j,2aj}(\tau)}
=
\begin{cases}
r_{j}(\tau),
& a\equiv\pm1\pmod5,\\[3pt]
-r_{j}(\tau)^{-1},
& a\equiv\pm2\pmod5.
\end{cases}
\end{equation}
Substituting \eqref{eq.expuniform} and \eqref{eq.quot} into
\eqref{eq.transformraw} yields \eqref{eq.gamact}.
\end{proof}

\end{document}
